\section{Conclusion}
\label{sec:conclusion}

This paper asked whether the structure of national mortgage markets shapes
the transmission of ECB monetary policy to house prices. Using a panel of
eleven euro area countries over the first twenty seven years of the euro, I
estimated local projections of real house prices on information cleaned
high frequency policy surprises and interacted the shock with a fixed
versus variable rate classification of mortgage markets. Theory predicts
larger price declines in variable rate countries after a tightening. The
data do not confirm this prediction. The estimated differential is small,
its sign depends on the choice of dynamic controls, and it is nowhere
robustly significant. What the data do show is a negative and significant
immediate house price response in the period since 2009, and a decisive
role for the information content of ECB announcements.

The main limitation is statistical power, and I want to state it
concretely. The design identifies the differential response from a single
euro area wide shock series interacted with a binary country trait. In
effect, eleven countries are split into two groups, and the interaction is
estimated from the difference between two group averages of responses to
the same 108 quarterly shocks. The natural consequence is wide confidence
intervals that contain both the theoretical prediction and its opposite. The binary classification adds to the
problem, because it discards the considerable variation in variable rate
shares within each group.

For policy the null result is still informative, within limits. At
the country level, the evidence provides no support for the view that ECB
decisions systematically split euro area housing markets along mortgage
system lines, but it cannot rule that view out either.

Future work could relax the constraints of this design in three ways. A
continuous variable rate share would use variation that the binary
indicator throws away. Regional house price data would multiply the cross
section and allow time fixed effects, which my design cannot include
because they would absorb the common shock, and \citet{battistini2025}
show that such data exist for the euro area. Household level data on
mortgage payment resets would test the budget channel directly instead of
inferring it from country groups. All three routes address the power
problem directly, with more variation instead of more horizons.
